\chapter{Dictionaries}

\section{Why This Matters}
In Machine Learning, dictionaries are the absolute backbone of data structure formatting.
\begin{itemize}
    \item \textbf{JSON Data:} APIs and datasets often return data in JSON format, which maps perfectly 1-to-1 with Python dictionaries.
    \item \textbf{Hyperparameters:} When training an ML model, you will often pass in a dictionary of configuration settings (e.g., \texttt{\{"learning\_rate": 0.01, "batch\_size": 32\}}).
    \item \textbf{Model Weights:} PyTorch and TensorFlow store neural network layers and weights inside dictionaries.
\end{itemize}

\section{Explanation}
A dictionary (\texttt{dict}) in Python is an unordered, \textbf{mutable} collection of \textbf{key-value pairs}. It is incredibly fast for looking up values based on a unique key. 

\textbf{Keys MUST be unique and immutable} (like strings, integers, or tuples).
\textbf{Values can be absolutely anything}.

\subsection{1. Syntax and Core Methods}
\begin{lstlisting}
student = {
    "name": "Alice",
    "age": 28,
    "courses": ["Math", "CompSci"]
}

# Modifying or Adding a new key-value pair
student["age"] = 29
\end{lstlisting}

\begin{itemize}
    \item \texttt{.get(key, default)}: Safely gets a value without crashing if it doesn't exist.
    \item \texttt{.keys()}: Returns all keys.
    \item \texttt{.values()}: Returns all values.
    \item \texttt{.items()}: Returns all \texttt{(key, value)} tuples.
    \item \texttt{.update(dict2)}: Merges another dictionary into this one.
\end{itemize}

\subsection{2. Safe Lookups}
Accessing a key that doesn't exist using brackets \texttt{[]} will crash your program with a \texttt{KeyError}. Always use \texttt{.get()} if you aren't 100\% sure the key exists.
\begin{lstlisting}
config = {"theme": "dark"}
# print(config["font_size"]) <--- KeyError! Program crashes.

font = config.get("font_size", 12) # Safe
\end{lstlisting}

\subsection{3. Iterating through a Dictionary}
To loop through a dictionary effectively, use the \texttt{.items()} method, which unpacks the key and value simultaneously.
\begin{lstlisting}
model_metrics = {"accuracy": 0.95, "loss": 0.02}

for metric_name, value in model_metrics.items():
    print(f"{metric_name}: {value}")
\end{lstlisting}

\subsection{4. Nested Dictionaries}
When working with JSON data, you will often deal with dictionaries inside dictionaries.
\begin{lstlisting}
dataset = {
    "user_101": {
        "name": "Bob",
        "purchases": [15.99, 20.00]
    }
}
print(dataset["user_101"]["purchases"][0]) # 15.99
\end{lstlisting}

\mlconnection{When you use Pandas, creating a DataFrame (a table of data) from scratch is often done using a dictionary where the keys are the column names and the values are lists representing the rows.}

\section{Solutions to Exercises}

\textbf{Exercise 1: Safe Access}
\begin{lstlisting}
user_profile = {"name": "Sarah", "email": "sarah@example.com"}
phone = user_profile.get("phone", "No phone provided")
\end{lstlisting}
\textit{Solution: \texttt{.get()} checks if the key "phone" exists. It safely returns the fallback string instead of crashing.}

\vspace{1em}
\textbf{Exercise 2: Dictionary Updating}
\begin{lstlisting}
default_config = {"lr": 0.01, "batch_size": 32, "epochs": 100}
user_config = {"batch_size": 64, "optimizer": "Adam"}
default_config.update(user_config)
\end{lstlisting}
\textit{Solution: The \texttt{.update()} method merges \texttt{user\_config} into \texttt{default\_config}, overwriting existing keys and adding new ones.}

\vspace{1em}
\textbf{Exercise 3: Nested Extraction}
\begin{lstlisting}
print(api_response["data"]["forecast"][1]["conditions"])
\end{lstlisting}
\textit{Solution: Chained indexing goes from outer dictionary \texttt{["data"]} to inner list \texttt{["forecast"]}, to index \texttt{[1]}, to key \texttt{["conditions"]}.}
